\section{Method}
\label{sec:method}

\begin{figure}[t]
    \centering
    \includegraphics[width=\linewidth]{figures/ufuzz.pdf}
    \vspace{-10pt}
    \caption{Overview of \system. A frozen structural index supplies certified
    mutation opportunities and stable canonical regions, while an isolated
    reference vault supplies evaluator-only answers and evidence. Each valid
    mutant changes exactly one input. Retrieval coverage and divergence guide
    search without access to failure labels or the offline region universe.}
    \label{fig:method}
    \vspace{-10pt}
\end{figure}

We develop \system, a fuzzing framework for discovering the memory-use
failures defined in \S\ref{sec:problem}. \system constructs seeds from memory
checkpoints and builds a frozen structural description of each checkpoint
(\S\ref{subsec:seed-init}). It then enumerates applicable mutation
obligations, generates and validates query or memory-state mutants
(\S\ref{subsec:mutation}), and executes valid mutants under a fixed probe
budget (\S\ref{subsec:execution}). Search uses only observable retrieval
behavior and non-oracle structural information. Evaluation references and
failure verdicts remain outside the search loop.

\subsection{Seed and Checkpoint Initialization}
\label{subsec:seed-init}

Let \(H_{1:t}=\langle h_1,\ldots,h_t\rangle\) denote the interaction history
up to time \(t\), and let \(M_t\) denote the resulting memory state. The
history and its stable benchmark source identifiers are retained as provenance
for \(M_t\). If a backend permits exact state copies, \system creates isolated
copies. Otherwise, it recreates a candidate checkpoint in a separate session
by replaying a frozen initialization artifact derived from \(H_{1:t}\);
equivalence is checked before paired fuzzing as described below. Thus, fuzzing
does not modify the original task state.

A seed is \(x=(M_t,q)\), where \(q\) is a query that may use information in
\(M_t\). Each original checkpoint \(M_i\) has a frozen Checkpoint Structural
Index
\[
\mathcal{I}(M_i)=(\mathcal{F}_i,\mathcal{Q}_i,\mathcal{P}_i,\mathcal{Z}_i),
\]
where \(\mathcal{F}_i\) contains certified canonical facts,
\(\mathcal{Q}_i\) contains structured query descriptions,
\(\mathcal{P}_i\) is the stable benchmark provenance registry, and
\(\mathcal{Z}_i\) contains structural certificates used for applicability and
validation. A canonical fact has the form
\[
f=(e,r,v,\tau,p),
\]
where \(e\), \(r\), \(v\), and \(\tau\) denote an entity, relation, value, and
temporal scope, and \(p\) is a stable benchmark source identifier. Backend
memory identifiers are adapter-local metadata rather than canonical fact
identities.

The index is partial but certified. Fields may come from benchmark structure,
deterministic or rule-based extraction, or LLM-assisted extraction, but an LLM
judgment alone cannot certify a field. Every field used for applicability must
be checked independently of mutant execution and traceable to stable
provenance. Information that cannot be certified remains unstructured; the
corresponding structured mutation opportunities are inapplicable.

For a structured query, we use
\[
\operatorname{QI}(q)=(e,s_e,r,s_r,\tau,c,\alpha),
\]
where \(s_e\) and \(s_r\) are the entity and relation spans, \(c\) contains
answer-affecting constraints, and \(\alpha\) is checkpoint-relative
answerability. This description contains no gold answer.

Separately, an evaluator-only ReferenceVault stores
\[
G(x)=(a,E^{+},E^{-}),
\]
where \(a\) is the expected answer, \(E^{+}\) is valid supporting evidence,
and \(E^{-}\) is evidence that must not support the current answer. We use
\(a=\bot\) for a checkpoint-relative unsupported query. Gold answers,
\(E^{+}\), \(E^{-}\), correctness annotations, oracle predicates, and failure
verdicts are never exposed to mutation selection, seed retention, scheduling,
or search-value computation. The structural index is non-oracle information;
the scheduler normally receives a frozen obligation projection produced from
it rather than enumerating the full index.

\paragraph{Shared initialization.}
For a fixed checkpoint \(M_i\), backend \(b\), and repetition seed \(s\), the
backend checkpoint is constructed once. Random Mutation, Unguided LLM
Mutation, and \system preferably start from exact clones of that state. If
cloning is unavailable, they replay the same frozen initialization artifact
with all controllable randomness, model versions, prompts, decoding settings, and
ingestion configuration fixed. Because identical replay input need not produce
identical backend state, \system performs an initialization-equivalence check
and records a stable state identity, fingerprint, or canonical state
projection before method-specific fuzzing. The exact fingerprinting procedure
is deferred to backend capability testing. A condition whose initialization
equivalence cannot be established is recorded as an initialization or
capability failure and excluded from the paired RQ1 comparison until resolved.
After equivalence is established, the original queries and parent observations
produce shared initial coverage \(C_0^{(i,b,s)}\), without a method index.

\subsection{Mutation}
\label{subsec:mutation}

Given a seed \(x=(M,q)\), \system changes one input at a time. A query mutation
produces \(x'=(M,q')\), whereas a memory-state mutation produces
\(x'=(M',q)\). An operator \(o\) defines a structural applicability condition,
a transformation, a pre-execution validator, and an evaluator-side relation
\(\mathcal{C}_o\). Retained mutants may become seeds in later rounds, but every
individual parent--mutant pair still changes exactly one input.

\paragraph{Query mutations.}
A query mutation keeps \(M\) fixed. We use three classes. A
\emph{meaning-preserving} mutation preserves the requested entity, relation,
temporal scope, answerability, and other answer-affecting constraints. Its
operators paraphrase, shorten removable context, add nonrestrictive context,
or change syntactic form. A \emph{target-changing} mutation changes exactly
one designated primary slot in \(\operatorname{QI}(q)\), chosen from entity,
relation, or temporal scope, while preserving the remaining slots and
constraints. The replacement must have provenance-backed canonical support
for the unchanged request structure. An \emph{unsupported} mutation uses a
structurally well-formed replacement for which support is absent from the
current checkpoint.

Unsupportedness requires a checkpoint-relative certificate
\(operatorname{Absent}(z,r,\tau,M)=1\), or the corresponding form without
\(\tau\). The verifier must inspect the full checkpoint provenance relevant to
the entity, relation, and scope; absence from a partial extracted fact table is
insufficient. Deterministic matching, structured lookup, and LLM-assisted
analysis may contribute, but an LLM assertion alone is insufficient. Failed
retrieval and incorrect response behavior are never used to establish
unsupportedness.

\paragraph{Memory-state mutations.}
A memory-state mutation keeps \(q\) fixed and applies a controlled operation
\(\delta\) to an isolated state, producing \(M'=\delta(M)\). \system uses the
backend's native memory operation or interaction semantics and validates the
observable transition before executing the mutant query.

We use three operations. An \emph{update} introduces a later, type-compatible
value for an existing entity--relation pair and remains within that pair's
canonical region. A \emph{deletion} removes or deactivates an existing fact.
It is applicable only if a public or native backend operation faithfully
realizes the designated change without modifying other relevant information.
For Graphiti, episode deletion is applicable only when deleting the episode
satisfies this condition; low-level edge deletion is not used to emulate it.
An \emph{unrelated change} modifies the value or state of an already indexed
canonical region that is structurally unrelated to
\(\operatorname{QI}(q)\). It cannot introduce a new canonical region.

The ReferenceVault constructs the corresponding \(G(x')\) separately for
offline evaluation. Mutation construction, validation, and scheduling do not
receive that reference.

\paragraph{Mutation obligations and applicability.}
At the beginning of round \(r\), \system freezes the seed corpus
\(\mathcal{S}_r\). A mutation obligation
\(\omega=(x,o,t)\) pairs a seed, an operator, and one concrete structurally
certified target. For backend \(b\), let
\[
\mathcal{T}_{o,b}(x)=
\left\{t:\operatorname{pre}^{\mathrm{struct}}_o(x,t)=1
\land \operatorname{faithful}_b(x,o,t)=1\right\}.
\]
The applicable obligation set is
\begin{equation}
\Omega_{r,b}=
\left\{(x,o,t)\mid x\in\mathcal{S}_r,
t\in\mathcal{T}_{o,b}(x)\right\}.
\label{eq:mutation-obligations}
\end{equation}
It is generated from the frozen structural index rather than manually
pre-written and remains fixed within the round. All compared methods use the
same backend-specific applicability mask.

An opportunity is \textsc{inapplicable} when the operator and target cannot
be used faithfully for the seed and backend. It is excluded from
\(\Omega_{r,b}\), is not generated, and does not enter the applicable
denominator. A candidate is \textsc{invalid} when its obligation was
applicable but the generated result fails validation. Its obligation remains
uncovered and may be retried under a fixed generation-attempt cap. The cap is
selected after the pilot and frozen before full evaluation.

Let \(\Omega_{r,b}^{\mathrm{exec}}\) contain obligations that have produced at
least one valid executed mutant. Diagnostic obligation coverage is
\begin{equation}
\mathrm{Cov}_{\mathrm{mut}}^{(r,b)}=
\frac{|\Omega_{r,b}^{\mathrm{exec}}|}{|\Omega_{r,b}|}.
\label{eq:mutation-coverage}
\end{equation}
Invalid candidates do not cover an obligation.

\paragraph{Validation.}
Every candidate is validated before execution using its QueryIntent,
structural certificates, stable provenance, and, for memory mutations,
observable operation or state information. Validation never uses the mutant
retrieval, mutant response, ReferenceVault, or failure oracle. An LLM may
propose a rewrite or contribute a rejection check, but its judgment alone
cannot accept a candidate. Detailed rules appear in
Appendix~\ref{app:mutation-validation}.

\subsection{Execution and Fuzzing}
\label{subsec:execution}

For a seed \(x\), let \(U(x)=(R_k(x),y(x))\) denote the observed ranked
retrieval and final response. Every retained seed stores the observation from
its first valid execution together with its exact state and query identity.
For initial seeds, these observations are produced during shared campaign
initialization. For a normal parent--mutant evaluation, \system reuses the
cached \(U(x)\) and executes only \(U(x')\).

A cached parent observation may be reused only when the exact state and query
identity are proven identical. If identity cannot be proven, reconstruction
or a parent execution is logged as overhead rather than silently reused or
silently charged as a mutant probe. The same reconstruction and accounting
protocol is used by all methods.

One campaign is the product of one original checkpoint \(M_i\), one backend
\(b\), one method \(h\), and one repetition seed \(s\). Its budget \(B\) counts
valid new mutant executions only. Shared initialization, invalid generation,
inapplicable opportunities, reconstruction, and separately required parent
executions do not consume \(B\); their costs are logged separately. Thus, a
fixed method/backend/repetition condition over \(N\) checkpoints performs
\(N B\) valid new mutant executions.

\paragraph{Failure detection.}
After execution, the evaluator may use \(G(x)\), \(G(x')\), retrieval
provenance, and responses to determine whether the required relation is
violated:
\begin{equation}
F_o(x,x')=
\mathbf{1}\!\left[
\neg\mathcal{C}_o\bigl(U(x),U(x');G(x),G(x')\bigr)
\right].
\label{eq:failure-detection}
\end{equation}
This evaluator runs outside fuzzing search. Its verdict cannot affect seed
retention, scheduling, retries, or search values. The executable definitions
of \(\mathcal{C}_o\) and answer matching, and the rule for deduplicating
distinct failures in \(\mathrm{Fail}@B\), remain to be specified before the
full evaluation.

\paragraph{Canonical regions.}
For a canonical fact with certified entity and relation, define
\[
\kappa(f)=(\texttt{entity-relation},\operatorname{norm}(e),
\operatorname{norm}(r)).
\]
For source unit \(p\), let \(K_i(p)\) contain the certified semantic keys of
facts whose provenance includes \(p\), and define
\begin{equation}
\operatorname{RegionsForSource}_i(p)=
\begin{cases}
K_i(p), & K_i(p)\neq\varnothing,\\
\{(\texttt{source},p)\}, & K_i(p)=\varnothing.
\end{cases}
\label{eq:source-regions}
\end{equation}
Thus, a source contributes a fallback region only when it yields no certified
entity--relation region. Residual unstructured text does not create a second
region for a source that already has a semantic key.

For original checkpoint \(M_i\), the frozen campaign universe is
\begin{equation}
U_i:=U(M_i)=
\bigcup_{p\in\mathcal{P}_i}\operatorname{RegionsForSource}_i(p).
\label{eq:region-universe}
\end{equation}
Descendant states do not define new universes. Update changes state within an
existing region, deletion may deactivate content while its region remains in
\(U_i\), and an unrelated change modifies another existing region. No valid
memory-state mutation creates a region outside \(U_i\).

For a retrieved backend entry \(m\), the adapter first constructs a certified
fact-level mapping
\[
\operatorname{FactMap}_{b,i}(m)=
\{f:m\text{ can reliably be shown to represent }f\},
\]
and defines the semantic region set
\[
R^{\mathrm{sem}}_{b,i}(m)=
\{\kappa(f):f\in\operatorname{FactMap}_{b,i}(m)\}.
\]
The separate conservative fallback set is
\[
R^{\mathrm{fallback}}_{b,i}(m)=
\{(\texttt{source},p):m\text{ has reliable provenance to }p
\land K_i(p)=\varnothing\}.
\]
We then define
\[
\rho_{b,i}(m)=R^{\mathrm{sem}}_{b,i}(m)\cup
R^{\mathrm{fallback}}_{b,i}(m)\subseteq U_i.
\]
When a source has no certified semantic key, reliable source provenance may
map the entry to its source fallback. When it has exactly one semantic key,
reliable source provenance suffices for that key. When it has several keys,
the entry maps only to facts that its content and provenance reliably
identify. An entry that cannot be disambiguated is marked ambiguous or
unmapped for coverage rather than assigned every key of its source. A
summarized entry may still map to several regions when it reliably represents
several certified facts. Mutation-generated values inherit their designated
target's key. Thus, coverage counts regions represented in the retrieved
entry, not every region present in its source unit. For ranked retrieval
\(R_k\), define
\[
\operatorname{Regions}_i(R_k)=
\bigcup_{m\in R_k}\rho_{b,i}(m).
\]

\paragraph{Retrieval-mapping quality.}
Because ambiguous and unmapped entries are conservatively excluded from
region coverage, we report
\[
\operatorname{MappingRate}=
\frac{|\{m:\rho_{b,i}(m)\neq\varnothing\}|}
{|\{m:m\text{ is a retrieved entry considered}\}|}.
\]
We also report mapped, ambiguous, and unmapped entry counts by
backend/benchmark condition and optionally by method as a sanity check. This
diagnostic does not guide scheduling, retention, or mutation selection; it
does not enter either part of the RegionCov fraction; and it is not a primary
RQ1 effectiveness metric.

\paragraph{Online search feedback.}
Let \(C_{t-1}^{\mathrm{online}}\) contain only regions observed during shared
initialization and earlier valid probes in the campaign. For a probe executed
on its actual state \(M_{i,t}\), let
\[
R_t=\operatorname{Regions}_i(R_k(q_t,M_{i,t})).
\]
Online coverage gain is
\begin{equation}
R_{\mathrm{cov},t}=
\frac{|R_t\setminus C_{t-1}^{\mathrm{online}}|}
{\max(1,|R_t|)},
\qquad
C_t^{\mathrm{online}}=C_{t-1}^{\mathrm{online}}\cup R_t.
\label{eq:online-region-gain}
\end{equation}
The scheduler may also use retrieval-pattern novelty and, for
meaning-preserving mutations, ranked retrieval divergence. It never receives
the complete \(U_i\), its denominator, unseen-region targets, or evaluator
information.

\paragraph{Offline region coverage.}
Let \(\mathcal{S}_0^{(i)}\) be the original benchmark query set for \(M_i\).
For backend \(b\) and repetition seed \(s\), shared initial coverage is
\begin{equation}
C_0^{(i,b,s)}=
\bigcup_{q\in\mathcal{S}_0^{(i)}}
\operatorname{Regions}_i(R_k(q,M_i)).
\label{eq:initial-region-coverage}
\end{equation}
These initialization executions do not consume \(B\). For method \(h\), let
\(x_j=(M_{i,j},q_j)\), \(1\leq j\leq B\), be its valid new mutant executions
in this checkpoint-level campaign. Then
\begin{equation}
C_B^{(i,b,h,s)}=C_0^{(i,b,s)}\cup
\bigcup_{j=1}^{B}
\operatorname{Regions}_i(R_k(q_j,M_{i,j})).
\label{eq:campaign-regions}
\end{equation}
Because all mappings are anchored, \(C_B^{(i,b,h,s)}\subseteq U_i\). We
define
\begin{equation}
\operatorname{RegionCov}_{i,b,h,s}@B=
\frac{|C_B^{(i,b,h,s)}|}{|U_i|},
\qquad
\Delta\operatorname{RegionCov}_{i,b,h,s}@B=
\frac{|C_B^{(i,b,h,s)}|-|C_0^{(i,b,s)}|}{|U_i|}.
\label{eq:region-coverage}
\end{equation}
For fixed \(b,h,s\), primary aggregation across \(N\) checkpoints is the
macro-average
\begin{equation}
\operatorname{RegionCov}_{b,h,s}@B=
\frac{1}{N}\sum_{i=1}^{N}\operatorname{RegionCov}_{i,b,h,s}@B,
\label{eq:macro-region-coverage}
\end{equation}
with the analogous definition for
\(\Delta\operatorname{RegionCov}@B\). A pooled result may be reported only as
a secondary diagnostic. The numerical budget, retrieval depth, retry cap,
repetition count, extraction model and version, normalization algorithm, and
manual-audit sample size are selected after the pilot and frozen before full
evaluation.
